\documentclass[10pt]{scrartcl}
% \usepackage[right]{lineno}
% \linenumbers

\raggedbottom
% \usepackage[scale=5]{draftwatermark}
% \usepackage{fullpage}
\usepackage[letterpaper,margin=1in]{geometry}
\geometry{
  letterpaper, 
  left=40mm,
  right=40mm,
}
\usepackage{fancyhdr}
\pagestyle{fancy}
% \rhead{
% Compiled: \today\ \currenttime
% }
% \lhead{\leftmark\rightmark}
% \rfoot{
% Compiled: \today\ \currenttime
% }

\usepackage{tikz-cd}

\usepackage{calc}
\usepackage[inline]{enumitem}
% \setlist[enumerate,1]{label=\Alph*.}

% make description label wrap
% from: https://tex.stackexchange.com/questions/320512/how-can-i-make-the-label-text-in-description-lists-wrap-in-the-label-area-over-m
% \makeatletter
% \def\enit@align@parright{%
%   \def\enit@align##1{%
%     \nobreak
%     \strut\smash{\parbox[t]\labelwidth{\raggedleft##1}}}}

% \def\enit@style@multilineright{%
%   \enit@align@parright
%   \enit@calcset\labelindent\z@{0pt}%
%   \enit@calcset\leftmargin\thr@@!}
% \makeatother

% \setlist[description]{labelwidth=\widthof{\bfseries 2022 Spring Exam 3},
%   style=multilineright}

\usepackage{graphicx}
\usepackage{float}
\usepackage{booktabs}
\usepackage{tcolorbox}
\usepackage{textcomp}
\usepackage{commath}
\usepackage{mathtools}
\usepackage{physics}
\usepackage{outlines}
\usepackage[colorlinks=true]{hyperref}

\usepackage{todonotes}

\usepackage{amsmath,amssymb}
\usepackage{amsthm}
\theoremstyle{definition}
\newtheorem{problem}{Problem}

\usepackage{siunitx}
\DeclareSIUnit{\feet}{ft}

\newcommand{\ZZ}{\mathbb{Z}}
\newcommand{\RR}{\mathbb{R}}
\newcommand{\vavg}{v_{\text{avg}}}
\newcommand{\inv}{^{-1}}

\everymath{\displaystyle}

\usepackage[yyyymmdd]{datetime}


\title{Calc 1 Exam Archive by Topic}
\date{
  2023 Oct 12 
% Last Compiled: \today\ \currenttime
}

\begin{document}
\maketitle
\tableofcontents

\part{Topics by exam (roughly)}
\begin{outline}
  \1 Before Calc 1
  \2 Algebra
  \2 Trig / Unit circle

  \1 Exam 1 topics
  \2 Basic derivative rules (product, quotient chain, trig)
  \2 Evaluating limits
  \3 Plug in
  \3 Algebraic manipulation
  \3 Special Trig limits
  \3 Squeeze Theorem
  \2 Limit definition of derivative
  \2 Continuity?
  \2 Intermediate
  Value Theorem (IVT)

  \1 Exam 2 topics
    \2 More derivative rules (log, exp, inverse trig)
    \2 Logarithmic differentiation
    \2 Implicit differentiation
    \2 Related rates
    \2 Linearization / Linear Approximation
    \2 Differentials
    \2 Mean value theorem?
    \2 Finding / identifying local and global extrema on a graph
    \2 The process of finding absolute min/max on a closed domain
    \2 First and second derivative tests
    \2 Calculating $\lim_{x\to\infty}$ and $\lim_{x\to-\infty}$
    \2 Curve sketching
    \2 Reading data off of a graph
      \3 Intervals where increasing/decreasing
      \3 Local mins and maxes
      \3 Intervals where concave up/down

  \1 Exam 3 topics
    \2 Applied optimization
      \3 Closed interval method
      \3 1st derivative test
      \3 2nd derivative test
    \2 L'H\^{o}pital's rule
    \2 Evaluating Integrals
      \3 Visually/geometrically (area under a curve)
      \3 Antiderivatives
      \3 U-substitution
      \3 Symmetry? (even/odd)
      \3 Initial value problem
    \2 Using integrals
      \3 ``Net change theorem''
      \3 Calculating displacement / distance given a velocity function and endpoints in time.
    \2 Approximating integrals
      \3 Left-Endpoint / Right-Endpoint / Midpoint
    \2 Fundamental Theorem of Calculus (FTC)

  \1 Exam Final
    \2 Area between two curves
    \2 Volumes of revolution
      \3 Disk / Washer method
      \3 Shell method
\end{outline}


\part{Past exam problems}

Within each grouping, problems are listed in reverse order by semester, in order by Exam type (1,2,3,final) and problem number.

\section{Limits}
I will sort the limit problems based on how they should be solved (just plug in, algebraic, trig limit, l'hopital).
This is usually not specified for problems. \\
Students should be able to determine the method to use without being told (or to experiment with different methods until one works) \\
Beware the indeterminate forms
\[
  \frac 00, \frac\infty\infty, \infty-\infty, 0\cdot\infty, 1^\infty, 0^0, \infty^0
\]
which are related as follows:
\begin{displaymath}
\begin{tikzcd}[column sep=.5in, row sep=.5in]
	{\dfrac00,\ \dfrac\infty\infty} &&& {\infty-\infty} \\%\arrow{lll}[swap]{\text{combine fractions}} \\
	0\cdot\infty \\
	{1^\infty,\ \infty^0,\ 0^0}
	\arrow["{\text{combine fractions}}"', from=1-4, to=1-1,anchor=center]
	\arrow["{\text{flip a factor into denominator}}"', from=2-1, to=1-1]
	\arrow["{\text{take ln()}}"', from=3-1, to=2-1]
\end{tikzcd}
\end{displaymath}

\subsection{Just plug in}
\begin{enumerate}
\item [2023 Spring Exam 1 \#7A] $\lim_{x\to2} \frac{5-x}{x^2}$
\item [2023 Spring Exam 1 \#7B] $\lim_{x\to1} (\ln(2-x) + 3x)$
\item [2022 Fall Exam 1 \#1A] $\lim_{x\to1} (x^6 + 4\sqrt x + 1)$
\item [2022 Fall Exam 1 \#1B] $\lim_{\theta\to\pi/2} \frac{\sin \theta}{\theta}$
\item [2022 Spring Exam 1 \#1A] $\lim_{x\to3} (5+x^2)$
\item [2022 Spring Exam 1 \#1B] $\lim_{\theta\to0}{\cos\theta}{1-\theta}$
\item [2019 Fall Exam 1 \#1A] $\lim_{\theta\to\pi} \frac{1-\sin\theta}{\theta}$
\item [2017 Spring Final] $\lim_{x\to0} \frac{e^x}{x+1}$
\item [2016 Spring Final] $\lim_{x\to\pi} \frac{\cos x}{x}$
\end{enumerate}

\subsection{Limit Law like (exam 1)}
These problems give the value of the limit at a point (e.g.~$\lim_{x\to2} f(x) =
5$) and ask for a limit of a more involved expression which uses the given
limit (e.g.~$\lim_{x\to2} 5(f(x))^2$).
Basically, want the student to use limit laws to just plug in the value.
\\
Won't copy problem. Just list
\begin{enumerate}
  \item [2023 Spring Exam 1 \#9 A,B]
  \item [2022 Fall Exam 1 \#3]
  \item [2022 Spring Exam 1 \#5]
  \item [2019 Fall Exam 1 \#6]
\end{enumerate}

\subsection{Algebraic technique (exam 1)}
Such as factoring, multiplying by conjugate, rewrite as a single fraction
\begin{enumerate}
\item [2023 Spring Exam 1 \#4A] $\lim_{t\to2}\frac{t^2-5t+6}{t-2}$
\item [2023 Spring Exam 1 \#4B] $\lim_{x\to0}\left( \frac1x - \frac{1}{x(x+1)} \right)$
\item [2022 Fall Exam 1 \#1C] $\lim_{t\to-1} \frac{t^2+3t+1}{t+1}$
\item [2022 Fall Exam 1 \#1D] $\lim_{x\to4} \frac{\sqrt x-2}{x-4}$
\item [2022 Spring Exam 1 \#7A] $\lim_{w\to3} \frac{w^2-9}{w-3}$
\item [2022 Spring Exam 1 \#7B] basically the same as 2022 Fall Exam 1 \#1D
\item[2021 Fall Final \#7A] $\lim_{x\to1} \frac{x^2+2x-3}{x^2+x-2}$
\item [2019 Fall Exam 1 \#1B] $\lim_{x\to-2} \frac{x^2-4}{x+2}$
\item [2019 Fall Exam 1 \#1C] $\lim_{x\to1} \frac{\sqrt t-1}{t-1}$
\item [2017 Fall Final] $\lim_{x\to2} \frac{x^3 - 4x}{2-x}$ \footnote{l'hopital can be used, but not needed}
\item [2016 Fall Final] $\lim_{x\to 3} \frac{x-3}{9x-x^3}$
\end{enumerate}

\subsection{Special Trig limits}
\begin{enumerate}
\item [2023 Spring Exam 1 \#7C] $\lim_{\theta\to0} \frac{3\sin \theta}{\theta}$
\item [2022 Spring Exam 1 \#1C] $\lim_{\theta\to0} \frac{7\sin \theta}{3 \theta}$
\item [2017 Fall Final] $\lim_{h\to 0} \frac{\sin(3h)}{\sin(2h)} + \frac{\cos(3h)}{\cos(2h)}$ \footnote{l'hopital can also be used for one of the terms}
\item [2017 Spring Final] $\lim_{\theta\to0} \frac{\sin(\theta^2)}{\theta}$ \footnote{l'hopital can also be used for this}
\end{enumerate}

\subsection{Squeeze Theorem (exam 1)}
Should be able to identify that these problems require squeeze theorem, without being prompted to use it.
Usually, problem will say something like ``mention any theorems used''.
\begin{enumerate}
\item [2023 Spring Exam 1 \#10] Find $\lim_{x\to0} m(x)$ provided that the function $m(x)$ satisfies $x+1 \le m(x) \le e^x$ for all $x \ne 0$.
\item [2022 Spring Exam 1 \#6] Find $\lim_{x\to2} m(x)$ provided that the function $m(x)$ satisfies $5x-4 \le m(x) \le x^2+x$ for all $x \ne 2$.
\item [2019 Fall Exam 1 \#2] Provided that $10x-25 \le h(x) \le x^2$ for all $x \ne 5$, find $\lim_{x\to5} h(x)$
\end{enumerate}

\subsection{Intermediate Value Theorem (IVT) (exam 1)}
For these problems, usually need to state that one is using IVT, and that the function is continuous, so that IVT applies.
\begin{enumerate}
\item [2023 Spring Exam 1 \#5] Show there is a root of $f(x) = 2\cos x - \sin x$ in the interval $(0, \frac\pi2)$.
\item [2022 Fall Exam 1 \#2] Show there is a root of $f(x) = x^3+2x-3$ in the interval $(0, 2)$.
\item [2022 Spring Exam 1 \#2] Show there is a root of $f(x) = e^x+x-2$ in the interval $(0, 2)$.
\end{enumerate}

\subsection{L'Hopital (exam 3 / final)}
Should show some work indicating that one has checked the limit is an indeterminate form, therefore allowing the use of l'Hopital
\begin{enumerate}
\item[2023 Spring Exam 3 \#5A] $\lim_{x\to\infty} \frac{3+5x}{e^x+7}$
\item[2023 Spring Exam 3 \#5B] $\lim_{\theta\to\pi} \frac{\sin\theta}{\theta-\pi}$
\item[2023 Spring Final \#1A] $\lim_{x\to\infty} \frac{-3x-9e^x}{7x+2e^x}$
\item [2022 Spring Exam 1 \#9A] Let $v(x) = \frac{2}{x}$. Find $v'(1)$
\item [2022 Fall Exam 3 \#3A] $\lim_{x\to\infty} \frac{x}{e^x}$
\item [2022 Fall Exam 3 \#3B] $\lim_{x\to-\infty} \frac{1+e^x}{\sqrt{e^x+1}}$
\item [2022 Fall Exam 3 \#3C] $\lim_{x\to0^{+}} x^x$
\item[2022 Spring Exam 3 \#7A] $\lim_{x\to\infty} \frac{e^x+2}{x^2+4}$
\item[2022 Spring Exam 3 \#7B] $\lim_{x\to1} \frac{\ln x}{x-1}$
\item[2022 Spring Final \#1A] $\lim_{x\to1} \frac{\ln x}{x-1}$
\item[2021 Fall Final \#7B] $\lim_{x\to0} \frac{\sin x+x^2}{\cos x-1+3x}$
\item [2019 Fall Exam 3 \#1A] $\lim_{\theta\to0} \frac{\theta^2+\theta^3}{\sin(\theta)}$ \footnote{This can be done using trig limit, but L'Hopital also works}
\item [2019 Fall Exam 3 \#1B] $\lim_{x\to\infty} \frac{3x+e^x}{5x^2+7}$
\item [2019 Fall Final \#1A] $\lim_{x\to\infty} \frac{3e^x + x}{7x-5e^x}$
\item [2018 Fall Exam 3] $\lim_{\theta \to 0} \frac{\sin(\theta^2)}{\theta^2}$
\item [2018 Fall Exam 3] $\lim_{x\to\infty} \frac{x \ln x}{x^2 + 3x}$
\item [2018 Fall Final] $\lim_{x\to1} \frac{\ln x}{x-1}$
\item [2017 Fall Exam 3] $\lim_{t\to0} \frac{1+t-\cos t}{t^2 + \sin(2t)}$
\item [2017 Fall Exam 3] $\lim_{x\to\infty} x^2 2^{-x}$
\item [2017 Fall Exam 3] $\lim_{x\to0^+} x^{2x}$
\item [2017 Fall Final] $\lim_{x\to\infty} (x^2+5)^{1/x}$
\item [2017 Spring Exam 3] $\lim_{x\to\infty} \frac{e^x + 5x}{x + 3}$
\item [2017 Spring Exam 3] $\lim_{\theta\to0} \frac{\sin(\theta^2)}{3\theta^2}$
\item [2016 Fall Exam 3] $\lim_{t\to 0} \frac{4t-\sin(2t)}{5t^2 + 3t}$
\item [2016 Fall Exam 3] $\lim_{x\to\infty} x \sin(\frac 2x)$
\item [2016 Fall Exam 3] $\lim_{x\to\infty} (5x)^{1/x}$
\item [2016 Fall Final] $\lim_{h\to 0}\frac{\tan(2h)}{\sin(5h)}$
\item [2016 Fall Final] $\lim_{x\to\infty} (5+x)^{1/x}$
\end{enumerate}

\subsection{End behavior / Horizontal Asymptotes (HA)}
\begin{enumerate}
\item[2023 Spring Exam 3 \#5C] $\lim_{x\to\infty} \frac{\sqrt{4x^2+3x}}{6x+7}$
\item[2022 Spring Exam 3 \#7C] $\lim_{x\to-\infty} \frac{-5x^4+x+2}{7x^4-x^3-2x+1}$
\item[2022 Spring Exam 3 \#7D] $\lim_{x\to\infty} \frac{\sqrt{x^2+4}}{2x+1}$
\item[2021 Fall Final \#7C] $\lim_{x\to-\infty} \frac{5x^4-x+1}{x^4-2x^2+3}$
\item [2019 Fall Exam 2 \#1C] $\lim_{x\to-\infty} \frac{2x^8-x^2+3}{7x^9+14x+6}$
\item [2018 Fall Exam 2] Find $\lim_{x\to-\infty} \frac{3x^9 - 7x + 3}{2 + 5x + 6x^9}$.
\item [2017 Fall Exam 2] Evaluate $\lim_{x\to-\infty} \frac{7-3x^5}{2x^5-15x^3}$.
\item [2016 Fall Exam 2] Evaluate $\lim_{x\to-\infty} \frac{x+3}{\sqrt{x^2-1}}$
\item [2016 Spring Final] $\lim_{x\to\infty} \frac{23 + x - 5x^5 - 3x^9}{4x^9 - 3x - 2}$
\item [2015 Spring Exam 1] Find horizontal asymptote(s) for $y = \frac{\sqrt{9x^2+5}}{2x-7}$
\end{enumerate}

\section{Limit definition of derivative}
For these problems, the limit definition \textbf{must} be used.
Both versions. $f'(a)$ or $f'(x)$.
\begin{enumerate}
\item [2023 Spring Exam 1 \#11] Let $v(x) = \sqrt{x}$. Using the limit definition of derivative, find $v'(4)$.
\item [2023 Spring Exam 2 \#3] Let $g(x) = 7x^2 + x$. Using the limit definition of derivative, find $g'(x)$.
\item [2023 Spring Final \#11] Find $f'(2)$ if $f(x) = x^2+5x$.
\item [2022 Fall Exam 1 \#5A] Let $v(x) = x^2+2x-1$. Find the slope of the tangent line to $v(x)$ at $x=0$
\item [2022 Spring Exam 2 \#3] $p(x) = x^2-3x+1$, Find $p'(x)$
\item [2022 Spring Exam 1 \#9] Let $v(x) = \frac2x$. Find $v'(1)$.
\item [2022 Spring Exam 2 \#3] $p(x) = x^2-3x+1$. Find $p'(x)$.
\item [2022 Spring Final \#11] $f'(3)$ if $f(x) = x^2-7$
\item [2021 Fall Final \#14] $f'(1)$ for $f(x) = x^2+5$
\item [2019 Fall Exam 1 \#5] $f(x) = \frac1x$. Find $f'(3)$.
\item [2019 Fall Final \#13] Find $f'(2)$ if $f(x) = x^2+3x$.
\item [2018 Fall Final] Find $f'(2)$ if $f(x) = x^2$.
\item [2017 Fall Final] $f'(x)$ for $f(x) = x^2 + x$.
\item [2017 Spring Final] Find $f'(2)$ if $f(x) = 3x^2$
\item [2016 Fall Final] Find $f'(x)$ for $f(x) = 3x^2 - x$
\item [2016 Spring Final] Find $f'(2)$ if $f(x) = x^2 + x$.
\end{enumerate}

\section{Taking Derivatives}

\subsection{Basic}
This includes all derivative rules (power, product, quotient, chain, exp, log, trig, inverse trig)


\begin{enumerate}
\item[2023 Spring Exam 2 \#1A] $\dod{}{x}\left( \frac{6}{x^2} + 6^x + e^2 \right)$
\item[2023 Spring Exam 2 \#1B] $\dod{}{x}\left( x^{5/2} \cdot \ln(x) \right)$
\item[2023 Spring Exam 2 \#1C] $\dod{}{t} \arctan(t^3-5t)$
\item[2023 Spring Exam 2 \#1D] $\dod{}{\theta} \sqrt{\cos(5\theta)}$
\item[2023 Spring Exam 2 \#1E] $\dod{}{x}\left( \frac{\tan(x) - x}{e^x+x} \right)$
\item[2023 Spring Final \#1D] $\dod{}{x} \left( \frac{\tan(x^2)}{\cos(x)+x^3} \right)$
\item[2023 Spring Final \#1E] $\dod{}{x}\left( \ln(x) \cdot \arctangent(x^3) \right)$
\item[2022 Fall Exam 1 \#7A] $\dod{}{x} (x^{10}+\sqrt{x})$
\item[2022 Fall Exam 1 \#7B] (Use Product Rule) $\dod{}{x} [(x^{1/3}+x^{-3})(x^2-1)]$
\item[2022 Fall Exam 1 \#7C] (Use Product Rule) $\dod{}{x} [(x^3-2)(x^\pi+x^{1/5})]$
\item[2022 Fall Exam 1 \#7D] $\dod{}{x} \left( \frac{x^2+6x}{\sqrt{x}} \right)$
\item[2022 Fall Exam 2 \#1A] $\dod{}{x}\left( 5e^x - \log_5 x + \arctan x \right)$
\item[2022 Fall Exam 2 \#1B] $\dod{}{x}\left( 20^x \cdot \cosine(x^{20}) \right)$
\item[2022 Fall Exam 2 \#1C] $\dod{}{x} \arccos(\sqrt{x})$
\item[2022 Fall Exam 2 \#1D] $\dod{}{\theta} \csc(\sine(\theta^2))$
\item[2022 Fall Exam 2 \#1E] $\dod{}{x} \left( \frac{e^{2x}+\ln(2x+1)}{x^6-7x} \right)$
\item[2022 Spring Exam 2 \#1A] $\dod{}{x}\left( \frac 3x + 2^x + \ln(3) \right)$
\item[2022 Spring Exam 2 \#1B] $\dod{}{x}\left( x^{3/2} \cdot \arctan x\right)$
\item[2022 Spring Exam 2 \#1C] $\dod{}{t} \tan(t^3-4t)$
\item[2022 Spring Exam 2 \#1D] $\dod{}{\theta} \ln(\sin(2\theta))$
\item[2022 Spring Exam 2 \#1E] $\dod{}{x} \left( \frac{\cos x- e^x}{e^x+2} \right)$
\item[2022 Spring Final \#1D] $\dod{}{x} \left( \frac{\cos(x^2)}{\sin x + e^x} \right)$
\item[2022 Spring Final \#1E] $\dod{}{x} \left( \sqrt{x^3+1} \cdot \arctan(x^2) \right)$
\item[2021 Fall Final \#3A] $\dod{}{x} \left( \arctan x + 2^x \ln x + \frac{1}{\sqrt[3]{x}} \right)$
\item[2021 Fall Final \#3B] $\dod{}{x} \left( \frac{\sin x}{x^4+2} \right)$
\item[2021 Fall Final \#4A] $\dod{}{x} \tan(x^3+5x+1)$
\item[2021 Fall Final \#4B] $\dod{}{x} \frac{1}{(x^2+3x+1)^5}$
\item[2021 Fall Final \#4C] $\dod{}{x} e^{\cos x}$
\item[2019 Fall Exam 1 \#11A] $\dod{}{x} \left( 5x^2+3x^{1/3}+3\pi \right)$
\item[2019 Fall Exam 1 \#11B] $\dod{}{x} \left( \frac{7}{\sqrt x}+2\sqrt x \right)$
\item[2019 Fall Exam 1 \#11C] $\dod{}{x} \left( \theta\cdot\tangent(\theta^3) \right)$
\item[2019 Fall Exam 1 \#11D] $\dod{}{x} \cos(\sine(\theta^2))$
\item[2019 Fall Exam 1 \#11E] $\dod{}{x} \left( \frac{x^{3/2}-5}{x^2+1} \right)$
\item[2019 Fall Exam 2 \#1A] $\dod{}{x} \left( \ln(x^3+7) \cdot \arctan(x^2+5) \right)$
\item[2019 Fall Exam 2 \#1B] $\dod{}{x} \left( \frac{2^x+5}{e^x+1} \right)$
\item [2019 Fall Final \#1D] $\dod{}{x} \left(\frac{\cos(x^7)}{\ln x+3}\right)$
\item [2019 Fall Final \#1E] $\dod{}{x} \left(\tangent(x^2)\cdot\arctan(x)\right)$
\item[2018 Fall Exam 2] $\dod{}{x}\left(x \cdot \arctangent(3x^2)\right)$
\item[2018 Fall Exam 2] $\dod{}{x}\left(x \cdot \frac{2^x - \ln x}{e^x + 1} \right)$
\item[2018 Fall Final] $\dod{}{x}\left( \frac{\tan x}{\ln x + 3} \right)$
\item[2018 Fall Final] $\dod{}{x} \sine(x^2) \cdot \arctangent(x)$
\item [2017 Fall Exam 2] $w(x) = \frac{\tan(e^x)}{1+x^3}$
\item [2017 Fall Exam 2] $h(x) = \sin\inv(\ln(x))$
\item [2017 Fall Final] $\dod{}{x} e^{3x} \tan\inv(x)$
\item [2017 Fall Final] $\dod{}{x} \frac{x+\tan x}{1-x^2}$
\item [2017 Summer Exam 2] $f(x) = \tan(\ln(6x^4+x^2))$
\item [2017 Summer Exam 2] $y = e^{\sin x}$
\item [2017 Spring Exam 2] $\dod{}{x}\left( 2^x - \frac{5}{x^2} + \ln(7) \right)$
\item [2017 Spring Exam 2] $\dod{}{x}\left(\sqrt x \cdot \tan x\right)$
\item [2017 Spring Exam 2] $\dod{}{x} \arctan(2t^2 - 3)$
\item [2017 Spring Exam 2] $\dod{}{x} \left( \frac{6\ln x-2x^3}{e^x+3} \right)$
\item [2017 Spring Final] $\dod{}{x} \frac{\tan x}{e^x + 5}$
\item [2017 Spring Final] $\dod{}{x} \cosine(x^3) \cdot \arctan(x)$
\item [2016 Fall Exam 2] $w(x) = x^3 e^{1/x}$
\item [2016 Fall Exam 2] $h(x) = \tan\inv(x^2)$
\item [2016 Fall Final] $f'(t)$ where $f(t) = \cos^2(2t+1)$
\item [2016 Fall Final] $\dod{}{x} x \ln(x^2 + 2)$
\item [2016 Fall Final] $\dod{}{x} \frac{e^{5x}}{x^2 + 1}$
\item [2016 Spring Exam 2] $\dod{}{x} \sin x \cdot 2^x$
\item [2016 Spring Exam 2] $\dod{}{x} \frac{\sqrt{x}}{\cos x}$
\item [2016 Spring Exam 2] $\dod{}{w} \arctan(5w^2 + 3)$
\item [2016 Spring Exam 2] $\dod{}{x}(e^{e^x})$
\item [2016 Spring Exam 2] $\dod{}{\theta} \tan\theta \cdot \ln\theta$
\item [2016 Spring Final] $\dod{}{x} \frac{e^x}{\sqrt x}$
\item [2016 Spring Final] $\dod{}{x} \ln(x) \cdot \sine(x^2)$
\end{enumerate}

\subsection{Implicit Diff}
Should know these use implicit diff.\ without being told.
Some of these ask for tangent line to the curve at a given point.

\begin{enumerate}
\item[2023 Spring Exam 2 \#4] Find $\dod{y}{x}$ if $x^4+xy+2y^3=7$.
\item[2023 Spring Exam 2 \#5] Find $\dod{y}{x}$ if $\sin(xy)=x^2$.
\item[2023 Spring Final \#2B] $\dod{y}{x}$ if $x^3+y^3 = 5-xy$
\item[2022 Fall Exam 2 \#3] Find $\dod{y}{x}$ if $\cosine(x^2y^3) = e^x$
\item[2022 Spring Exam 2 \#4] Find $\dod{y}{x}$ if $x^2y - x = e^y + 1$.
\item[2022 Spring Final \#2B] $\dod yx$ if $y^4+xy = x^3-x+2$
\item[2021 Fall Final \#2] $\dod yx$ if $y^5x^3 = 4x+y+10$.
\item[2019 Fall Exam 2 \#4] $\sin(xy) = x^2$. Compute $\dod yx$.
\item[2019 Fall Final \#2B] $\dod yx$ if $y^4 + xy = x^3-x+2$.
\item[2018 Fall Exam 2] Compute $\dod{y}{x}$ for $x^2y - e^y = x + 1$
\item[2018 Fall Final] Find $\dod{y}{x}$ for $x^3 + xy + y^4 = 5$
\item[2017 Fall Exam 2] Find an equation for the tangent line to the curve $x^2 y + y^2 = x^3 - 3$ at the point $(2, 1)$.
\item[2017 Fall Final] Find the equation of the tangent line to the curve $xy + y^2 = 2x - 1$ at $(2, 1)$.
\item[2017 Summer Exam 2] Compute $\dod{y}{x}$ for $3xy^2 = y^3 - \cos x$
\item[2017 Spring Exam 2] Find $\dod{y}{x}$ if $\sin(xy^2) = x^2$.
\item[2017 Spring Final] Find $\dod yx$ for $x^2 - xy + y^3 = 5$
\item[2016 Spring Exam 2] Find $\dod yx$ if $x^2y^3 = e^x - y^2$.
\item[2016 Fall Exam 2] Find an equation for the tangent line to the curve $xy + 7 = x^3 + y^3$ at the point $(2, 1)$.
\item[??] Find an equation of the tangent line to the curve $x^2 y^3 - x^3 y^2 = 4$ at the point $(1, 2)$.
\item[2016 Fall Final] Find the equation of the tangent line to the curve $x^3 + y^2 = 5y + 4$ at $(2, 1)$.
\item[2016 Spring Final] Find $\dod yx$ for $x^3 + y^3 = 5xy$
\end{enumerate}

\subsection{Logarithmic Diff}
Should be able to know these use Log.\ Diff.\ without being told.

\begin{enumerate}
\item[2023 Spring Exam 2 \#8] Find the derivative of $f(x) = x^{\cos(x)}$
\item[2023 Spring Final \#2C] $k'(x)$ if $k(x) = x^{\sin(x)}$
\item[2022 Fall Exam 2 \#2] Find the derivative of $h(x) = 7^{x+5} \cdot x^{7\tan x}$
\item[2022 Spring Exam 2 \#8] Find the derivative of $h(x) = x^{8x}$
\item[2022 Spring Final \#2C] $k'(x)$ if $k(x) = x^{3x}$
\item[2021 Fall Final \#13] Find $\dod yx$ for $y=\frac{(x-3)^x}{x^4}$
\item[2019 Fall Exam 2 \#6] Same as 2023 Spring Exam 2 \#8
\item[2019 Fall Final \#2C] $k'(x)$ if $k(x) = \frac{x^x}{\sin[7](x)}$
\item[2018 Fall Exam 2] Find the derivative of $h(x) = x^{5x^2}$
\item[2018 Fall Final] Find the derivative of $f(x) = x^{3x}$
\item[2017 Fall Exam 2] $y = x^{5x}(2 + 3x^2)^4$
\item[2017 Summer Exam 2] Find the derivative of $y = \frac{(x+2)^2}{(x+5)(3x-4)}$. \footnote{can be done without log.\ diff}
\item[2017 Spring Exam 2] Find the derivative of $h(x) = e^x \cdot x^{5\cos(x)}$
\item[2017 Spring Final] Find $\dod yx$ for $y = x^{7x}$
\item[2016 Spring Exam 2] Find the derivative of $h(x) = \frac{x^x}{(3x^2+4)^5}$
\item[2016 Spring Exam 2] Find $\dod yx$ if $y = x^{2x}(1-x)^7$.
\item[2016 Spring Final] Find $\dod yx$ for $y=x^{\cos x}$
\end{enumerate}

\subsection{Computing derivatives without the functions (given graph or table)}
For these problem, I won't be reproducing the graph/table, so I will just list the location of such exercises.
\begin{enumerate}
\item [2023 Spring Exam 2 \#7]
\item [2022 Spring Exam 2 \#7]
\end{enumerate}

\section{Equation of tangent line}
Some of these problems want the derivative to be computed using the limit definition first.
This will be ignored (maybe incorporated into the limit definition section).
\begin{enumerate}
\item [2023 Spring Exam 1 \#11B] Let $v(x) = \sqrt{x}$. Find the equation of the tangent line to $y=v(x)$ at $x=4$.
\item [2023 Spring Exam 2 \#2] $y=\sin x + 2$ at $x=0$
\item [2022 Fall Exam 1 \# 5B] Let $v(x) = x^2+2x-1$. Find the equation of the tangent line to $y=v(x)$ at $x=0$.
\item [2022 Spring Exam 1 \#9A] Find $v'(1)$.
\item [2022 Spring Exam 2 \#2] $y = \sqrt x$ at $x=4$
\item [2021 Fall Final \#1] $y=x^5+3x$ at $x=1$
\item [2019 Fall Exam 1 \#10] Given that $v(5) = 3$ and $v'(5) = -2$, find an equation of the tangent line to the graph of $y = v(x)$ at $x=5$.
\end{enumerate}

\section{Related rates}
\textbf{\color{red}{Always include units!}}

\begin{enumerate}
\item[2023 Spring Exam 3 \#6] A 5 meter ladder is leaning against a wall. If the top of the ladder slides down the wall at a rate of $\frac12$ m/s, how fast will the bottom of the ladder be moving along the ground when the top of the ladder is 4 meters above the ground?
\item[2023 Spring Final \#14] Boyle's Law states that when a sample of gas is compressed at a constant temperature, the pressure $P$ and the volume $V$ satisfy the equation $PV = C$, where $C$ is a constant. Suppose that at a certain instant, the volume is $\SI{40}{cm^3}$, the pressure is 100 kPa, and the pressure is increasing at a rate of 20 kPa/min. At what rate is the volume changing at this instant?
\item[2022 Fall Exam 2 \#6] Two airplanes are flying in the air at the same height: airplane A is flying east at 100 mi/h and airplane B is flying north at 200 mi/h. If they are both heading to the same airport, located 30 miles east of airplane A and 40 miles north of airplane B, at what rate is the distance between the airplanes changing?
\item[2022 Spring Exam 2 \#5] A hot air balloon rising vertically is tracked by an observer located 1 mile from the lift-off point. At a certain moment, the angle between the observer's line of sight and the horizontal is $\tfrac\pi4$, and it is changing at a rate of $\tfrac{1}{8}$ rad/min. How fast is the balloon rising at this moment?
\item[2022 Spring Exam 2 \#9] The volume of a sphere of radius $r$ is $V = \frac43 \pi r^3$. If the radius of a sphere is 2 cm and increasing at a rate of $\tfrac12$ cm/sec, find the rate at which the volume of the sphere is increasing.
\item[2022 Spring Final \#14] The base of a triangle is increasing at a rate of 2 ft/s, and the height of the triangle is increasing at a rate of 3 ft/s. Find the rate at which the area of the triangle is changing when the base is 4 ft and the height is 7 ft.
\item[2021 Fall Final \#8] The radius of a circle increases at a rate of 2 m/sec. Find the rate at which the area of the circle increases when the radius is 5 m.
\item[2019 Fall Exam 2 \#3] A 5 foot ladder is leaning against a wall. If the top of the ladder slides down the wall at a rate of 2 ft/s, how fast will the bottom of the ladder be moving away from the wall when the top is 3 feet above the ground?
\item[2019 Fall Final \#10] Boyle's Law states that when a sample of gas is compressed at a constant temperature, the pressure $P$ and the volume $V$ satisfy the equation $PV = C$, where $C$ is a constant. Suppose that at a certain instant, the volume is $\SI{30}{cm^3}$, the pressure is 100 kPa, and the pressure is increasing at a rate of 20 kPa/min. At what rate is the volume changing at this instant?
\item[2018 Fall Exam 2] A hot air balloon rising vertically is tracked by an observer located 2 miles from the lift-off point. At a certain moment, the angle between the observer's line of sight and the horizontal is $\tfrac\pi4$, and it is changing at a rate of $\tfrac{1}{10}$ radians/minute. How fast is the balloon rising at this moment?
\item[2018 Fall Final] A 5-foot ladder rests against a wall. The bottom of the ladder slides away from the wall at a rate of \SI{2}{ft/s}. How fast is the top of the ladder sliding down the wall when the bottom of the ladder is \SI{3}{ft} from the wall?
\item[2017 Fall Exam 2] A rocket is launched vertically upward from a point $P$. At the same time, a car is driving on a straight line away from the point $P$. Use related rates to determine the rate that the distance between the rocket and the car is increasing at the instant when the rocket is 3 miles up and traveling 500 miles per hour, and the car is 4 miles away from $P$ and traveling 30 miles per hour. (Assume the ground is flat.)
\item[2017 Fall Final] Consider a rectangle with edges of length $x, y$. If $x$ is increasing at a rate of \SI{5}{m/sec} and $y$ is decreasing at a rate of \SI{2}{m/sec}, at what rate is the area $A$ of the rectangle changing when $x = \SI{3}{m}$ and $y = \SI{4}{m}$?
\item[2017 Summer Exam 2] A plane is flying away from you at 500 mph at a height of 3 miles. How fast is the plane's distance from you increasing at the moment when the plane is flying over a point on the ground 4 miles from you?
\item[2017 Spring Exam 2] Boyle's Law says that when a sample of gas is compressed at a constant temperature, the pressure $P$ and the volume $V$ satisfy the equation $PV = C$, where $C$ is a constant. Suppose that at a certain instant, the volume is $\SI{300}{cm^3}$, the pressure is 100 kPa, and the pressure is increasing at a rate of 20 kPa/min. At what rate is the volume changing at this instant?
\item[2017 Spring Exam 2] The length of a rectangle is increasing at a rate of 2 ft/s, and its width is increasing at a rate of 3 ft/s. At what rate is the area of the rectangle increasing when the length is 6 ft and the width is 7 ft?
\item[2017 Spring Final] Suppose that an oil spill from a ruptured tanker spreads in a circular pattern. If the radius of the oil spill increases at a constant rate of \SI{2}{ft/sec}, how fast is the area of the spill increasing when the radius is \SI{10}{ft}?
\item[2016 Fall Final] Consider a right triangle with edges of length $x, y, z$, with $z$ the hypotenuse. If $x$ is increasing at a rate of \SI{5}{m/sec} and $z$ is increasing at a rate of \SI{7}{m/sec}, at what rate is $y$ increasing when $x = \SI{3}{m}$ and $z = \SI{5}{m}$?
\item[2016 Spring Final] A hot air balloon rising vertically is tracked by an observer located 5 miles from the lift-off point. At a certain moment, the angle between the observer's line of sight and the horizontal is $\frac \pi 4$, and it is changing at a rate of $\frac{1}{10}$ radians/minute. How fast is the balloon rising at this moment?
\end{enumerate}

\section{Linear Approximation/Linearization}
\begin{enumerate}
\item [2023 Spring Exam 3 \#7] Find the linearization of $g(x) = \ln x$ at $x=1$. Use this to estimate $\ln(1.07)$.
\item [2023 Spring Final \#10] Use the linearization of $u(x) = \sqrt x$ at $x=25$ to approximate $\sqrt{24}$.
\item [2022 Fall Exam 2 \#4] Let $f(x) = \frac{\sin x}{x-\pi+1}$. Find the linearization for $f(x)$ at $x=\pi$. Use the linearization to approximate $f(\pi+0.02)$
\item [2022 Spring Exam 3 \#6] Find the linearization of $g(x) = \sin x$ at $x=0$. Use it to estimate $\sin(0.01)$.
\item [2022 Spring Final \#10] Let $p(t)$ denote the position of a particle in cm after $t$ seconds. Suppose that we are unable to write down a nice formula for $p(t)$, but we happen to know that when $t=2$ seconds, the particle has position 7 cm and velocity 30 cm/s. Find the linearization of $p(t)$ at $t=2$ seconds, and use it to approximate the location of the particle when $t=1.9$ seconds.
\item [2019 Fall Exam 2 \#2] Let $s(t)$ denote the position of a particle in cm after $t$ seconds. Suppose that we are unable to write down a nice formula for $s(t)$, but we happen to know that when $t=1$ second, the particle has position 5 cm and velocity 2 cm/s. Find the linearization of $s(t)$ at $t=1$ seconds, and use it to approximate the location of the particle when $t=0.9$ seconds.
\item [2019 Fall Final \#6] Use a linearization of $u(x) = \sqrt{x}$ at $x=25$ to approximate $\sqrt{26}$.
\item [2018 Fall Exam 2] Same as 2019 Fall Final \#6 %Find the linearization of $g(x) = \sqrt{x}$ at $x = 25$. Use this to estimate $\sqrt{26}$.
\item [2018 Fall Final] Use the linearization of $u(x) = \ln x$ at $x = 1$ to approximate $\ln(0.9)$.
\item [2017 Fall Exam 2] Find the linear approximation of $f(x) = \cos(x)$ near $x = \frac{\pi}{4}$. Use this to estimate $\cos(\frac\pi4+\frac{1}{10})$.
\item [2017 Fall Final] Find the linear approximation of $f(x) = \sqrt{x}$ near $x = 9$. Use this to estimate $\sqrt{8.9}$.
\item [2017 Spring Exam 3] Find linearization for $g(x) = \ln(x)$ at $x = 1$. Use it to approximate $\ln(1.15)$.
\item [2017 Spring Final] Use a linearization of $u(x) = \sqrt{x}$ at $x = 9$ to approximate $\sqrt{9.6}$.
\item [2016 Fall Exam 2] Find the linear approximation of $f(x) = \sqrt[3]{x}$ near $x = 8$. Use it to estimate $\sqrt[3]{8.1}$.
\item [2016 Fall Final] Find the linear approximation of $f(x) = \sqrt{x}$ near $x = 4$. Use this to estimate $\sqrt{4.1}$.
\item [2016 Spring Exam 3] Same as 2017 Spring Final %Find the linearization of $w(x) = \sqrt{x}$ at $x = 9$. Use it to estimate $\sqrt{9.6}$.
\item [2016 Spring Final] Use linearization to approximate $\sin(0.01)$.
\item [2015 Fall Exam 2] Find an equation of the tangent line to the curve $y = \sqrt[5]{x}$ at $x = 32$, and use this to approximate $\sqrt[5]{33}$.
\end{enumerate}

\newpage
\section{Differentials}
\begin{enumerate}
\item [2023 Spring Exam 3 \#8] The surface area $A$ of a sphere of radius $r$ is given by $A = 4\pi r^2$. Find the differential $\dif A$.
\item [2022 Spring Exam 3 \#4] Find the differential $dy$ if $y = \cosine(x^2+3)$.
\item [2019 Fall Exam 2 \#9] The volume of a sphere is $V = \frac43 \pi r^3$, where $r$ denotes the radius. Find the differential $\dif V$ in terms of $r$ and $\dif r$.
\item [2018 Fall Exam 2] Let $V$ denote the volume of a cube of side length $x$. Find the differential $dV$ in terms of $x$ and $dx$.
\item [2017 Fall Exam 2] The volume of a cone of height 9 ft is given by $V = 3\pi r^2$. Estimate the change in volume using the differentials $dV$ and $dr$, if $r = 5$ ft and increased by $\frac{1}{10}$ ft.
\item [2017 Spring Exam 3] The volume of a sphere is given by $V = \frac 43 \pi r^3$. Find the differential $dV$.
\item [2016 Fall Exam 2] The volume of a sphere is given by $V = \frac 43 \pi r^3$. Estimate the change in volume by calculating $dV$, given that $r = 3$ inches and $dr = \frac{1}{12\pi}$ inches.
\item [2016 Spring Exam 3] Find $dy$ if $y = \cosine(4x^2)$.
\end{enumerate}

\newpage
\section{Abs min/max over closed interval (Closed Interval Method)}
\begin{enumerate}
\item [2023 Spring Exam 3 \#2] $w(x) = x^3-3x+2$ on $[0,2]$
\item [2023 Spring Final \#13] Same as 2023 Spring Exam 3 \#2
\item [2022 Fall Exam 2 \#7] Let $f(x) = x^4-2x^2$. Find abs min/max of $f(x)$ on $[-2,2]$.
\item [2022 Spring Exam 3 \#5] $w(x) = x^3-3x^2+1$ on $[-1,1]$
\item [2022 Spring Final \#13] $w(x) = (x-1)e^x$ on the interval $[-1,1]$.
\item [2019 Fall Exam 2 \#5] Same as 2023 Spring Exam 3 \#2
\item [2019 Fall Final \#7] $w(x) = x + \cos(x)$ on $[0,\pi]$.
\item [2018 Fall Exam 2] $w(x) = x + \sin x$ on $[0, 2\pi]$.
\item [2018 Fall Final] $w(x) = x - \sqrt{x}$ on $[0, 4]$.
\item [2017 Fall Exam 2] Determine the absolute minimum and absolute maximum value of the function $f(x) = x^4 - 2x^2$ over the interval $[0, 2]$.
\item [2017 Spring Exam 3] Find the absolute min and max of $w(x) = 2x^3 - 9x^2 + 3$ on $[-1, 1]$.
\item [2017 Spring Final] $v(x) = x^3 + 3x^2 + 1$ on $[-1, 1]$
\item [2016 Spring Exam 3] Find absolute min/max of $g(x) = x^3 - 3x^2 + 4$ on $[-1, 1]$.
\item [2016 Spring Final] $v(x) = x^3 - 3x + 1$ on $[0, 2]$
\item [2016 Fall Exam 2] Find absolute min/max value of $f(x) = x^3 - 3x$ over the interval $[0, 2]$.
\item [2015 Fall Exam 2] Find the $x$ and $y$ coordinates of the absolute maximum of the function $y = x^3 - x^2 - x$ on the closed interval $0 \le x \le 2$
\end{enumerate}


\section{Read off increasing/decreasing, local min/max, CU/CD, inflection pt from graph of derivative}
Won't duplicate problem. Just list place.
\begin{enumerate}
\item [2023 Spring Exam 3 \#3]
\item [2023 Spring Final \#6]
\item [2022 Spring Exam 3 \#2]
\item [2022 Spring Final \#6]
\item [2019 Fall Final \#9]
\item [2019 Fall Exam 2 \#7]
\end{enumerate}


\section{Easy interpretation of local min/max (application of 1st derivative test)}
Won't duplicate problem. Just list place.
\begin{enumerate}
\item [2023 Spring Exam 3 \#4]
\item [2022 Fall Exam 3 \#2] This one looks harder than usual.
\item [2022 Spring Exam 3 \#3]
\item [2019 Fall Exam 2 \#8]
\end{enumerate}

\newpage
\section{Curve sketch process}
\begin{tcolorbox}
  By ``the full process for curve sketching'' I mean:
  \begin{enumerate}
  \item Finding intercepts
  \item Finding asymptotes (HA, VA)
  \item Finding critical numbers and sign chart for $f'$
  \item Finding inflection points and sign chart for $f''$
  \item Sketching the curve using the above data.
  \end{enumerate}
\end{tcolorbox}
Won't fully duplicate problem. Just have functions.
The problem is a bunch of parts, walking through the curve sketch process (intercepts, HA, VA, increasing, decreasing, local min/max, concave up/down, inflection points) allowing one to then accurately sketch a graph of the function.

\begin{enumerate}
\item [2023 Spring Exam 3 \#1]
$\begin{aligned}
    f(x) = \frac{x^2}{x^2+3}
    \qquad
    f'(x) = \frac{6x}{(x^2+3)^2}
    \qquad
    f''(x) = \frac{-18(x-1)(x+1)}{(x^2+3)^3}
  \end{aligned}$
\item [2022 Fall Exam 3 \#1]
  $
    f(x) = \sqrt{x^2+2x}
    \qquad
    f'(x) = \frac{x+1}{\sqrt{x^2+2x}}
    \qquad
    f''(x) = \frac{-1}{(x^2+2x)^{3/2}}
  $
\item [2022 Spring Exam 3 \#1]
  $
    f(x) = x^2(x-3)
    \qquad
    f'(x) = 3x(x-2)
    \qquad
    f''(x) = 6(x-1)
  $
\item [2021 Fall Final \#9]
  $
    f(x) = \frac{x}{(x^2+9)^2}
    \qquad
    f'(x) = \frac{3(3-x^2)}{(x^2+9)^3}
    \qquad
    f''(x) = \frac{12x(x^2-9)}{(x^2+9)^4}
  $
\item [2019 Fall Exam 2 \#10]
  $
    f(x) = \frac{x^2-1}{x^2+3}
    \qquad
    f'(x) = \frac{8x}{(x^2+3)^2}
    \qquad
    f''(x) = \frac{-24(x^2-1)}{(x^2+3)^3}
  $
\item [2018 Fall Exam 2] Sketch $f(x) = x^2(x+3)$, which has derivatives $f'(x) = 3x(x+2)$, $f''(x)=6(x+1)$. Do the entire process.
\item [2017 Fall Exam 2] For $g(x) = 20x^3 - 3x^5$, draw the number line for $g'(x)$ and $g''(x)$. Classify the critical points.
\item [2017 Fall Exam 2] Let $f(x) = \frac{2-x}{x^2 - 1}$. Then $f'(x) = \frac{x^2- 4x+1}{(x^2-1)^2}$. Find asymptotes, intercepts, critical points. Classify the critical points.
\item [2017 Fall Final] Let $f(x)$ be a function with $f'(x) = x^2(x^2-4)(x-7)$. Find the critical points of $f(x)$, and the intervals where $f(x)$ is increasing / decreasing. Classify each critical point as a local min, local max, or neither.
\item [2017 Fall Final] Let $g(x) = 2x^6 - 5x^4$. Determine the intervals where $g(x)$ is concave up / concave down. Determine all inflection points.
\item [2017 Summer Exam 2] Let $f(x) = x^3 - 2x^2 + x - 1$. Find when $f$ is concave up/down, increasing/decreasing, the inflection points, the critical points, and the local minima and maxima.
\item [2017 Spring Exam 3] Sketch $f(x) = \frac{x^2}{x^2 + 3}$. It has $f'(x) = \frac{6x}{(x^2+3)^2}$ and $f''(x) = \frac{-18(x^2-1)}{(x^2+3)^3}$. Do the entire process.
\item [2017 Spring Final] Given $w''(x) = \frac{3-x}{x^2 + 7}$, find the intervals where $w(x)$ is concave up / concave down. Determine inflection points.
\item [2016 Fall Exam 2] Let $g(x) = x^5 - \frac{20}{3} x^3$. Draw the number line for $g'$ and $g''$. Classify the critical points. State when $g(x)$ is concave up.
\item [2016 Fall Exam 2] Let $f(x) = \frac{2x^2-6x}{4-x^2}$. Then $f'(x) = \frac{-2(3x^2-8x+12)}{(4-x^2)^2}$. Sketch $f(x)$ (skip using $f''$ or inflection points).
\item [2016 Fall Final] Let $g(x) = 3x^5 + 20x^3$. Determine intervals where $g(x)$ is concave up / concave down. Determine all inflection points.
\item [2016 Fall Final] Let $f(x) = x^2 (x-4)^3$. Given $f'(x) = x(x-4)^2(5x-8)$. Find the critical points of $f(x)$. Find the intervals where $f(x)$ is increasing / decreasing, and classify the critical points (local min / local max / neither).
\item [2016 Spring Exam 3] Sketch $f(x) = \frac{x^2- 1}{x^2 + 3}$. It has $f'(x) = \frac{8x}{(x^2+3)^2}$ and $f''(x) = \frac{-24(x^2-1)}{(x^2+3)^3}$. Do the entire process.
\item [2016 Spring Final] Given $w''(x) = \frac{2(x-1)}{x^2 + 3}$, find the intervals where $w(x)$ is concave up / concave down. Determine inflection points.
\item [2015 Fall Exam 3] Let $f(x) = \left( \frac{x-2}{x} \right)^2$. Then $f'(x) = \frac{4(x-2)}{x^3}$ and $f''(x) = \frac{-8(x-3)}{x^4}$. Do the entire process.
\end{enumerate}

\newpage
\section{Applied Optimization}
In these problems, one must justify why the critical point found is the absolute min/max for the implicit domain of the function being optimized.\\
Some of these problems come with diagram, which I won't reproduce. \\
\textcolor{red}{Include units in final answer!}

\begin{enumerate}
\item [2023 Spring Exam 3 \#9] Suppose that you have 24 meters of fencing to make two adjacent rectangular kennels of length $x$ meters and width $y$ meters. Find the values of $x$ and $y$ that maximize the enclosed area.
\item [2023 Spring Final \#7] Let $p(x) = 75-x^2$ be the price in dollars per meal that a chef can charge if they sell $x$ meals. Revenue is the total amount of money received from the sale of $x$ meals. Find the meal price that will maximize revenue.
\item [2022 Fall Exam 3 \#5] A rectangular open-topped aquarium is to have a square base and volume $\SI{9}{m^3}$. The material for the base costs \$2 per $\unit{m^2}$, and the material for the sides costs \$3 per $\unit{m^2}$. What dimensions minimize the cost of the aquarium?
\item [2022 Spring Exam 3 \#8] Suppose that when your bakery sells cakes for $x$ dollars each, your total profit is $P(x) = -x^2+20x-75$ dollars. In order to maximize total profit, how much should your bakery charge per cake?
\item [2022 Spring Exam 3 \#9] A homeowner with 20 feet of fencing wants to enclose a rectangular area against the side of her house. What dimensions will maximize the fenced-in area? (Note that three sides of the rectangle will be formed from fencing, and the house will serve as the fourth side of the rectangle.)
\item [2022 Spring Final \#7] When a company charges $x$ dollars per backpack, it makes a total profit of $P(x) = -2x^2+200x-50$ dollars. If the company wants to maximize total profit, what should it charge per backpack?
\item [2019 Fall Exam 3 \#8] A rectangular open-topped box is to have volume 18 \unit{m^3}. The length of the base is to be three times its width. What dimensions minimize the amount of material needed to make the box?
\item [2019 Fall Final \#14] A rectangular open-topped box is to have a square base and volume 8 \unit{ft^3}. If material for the base costs \$2 per \si{ft^2} and the material for the sides costs \$1 per \si{ft^2}, what dimensions minimize the cost of the box?
\item [2018 Fall Exam 3] A rectangular open-topped aquarium is to have a square base and volume \SI{8}{m^3}. The material for the base costs \$2 per \si{m^2}, and the material for the sides costs \$1 per \si{m^2}. What dimensions minimize the cost of the aquarium?
\item [2018 Fall Final] When a company charges $x$ dollars per chair, it makes a total profit $P(x) = -2x^2 + 200x - 50$ dollars. If the company wants to maximize total profit, what should it charge per chair?
\item [2017 Fall Exam 3] A box with square base and open top is formed from two materials. The base costs \$4 per square foot, while the four sides cost \$1 per square foot. If the total cost for the base and four sides is fixed to ber \$120, find the dimensions that maximize the volume of the box.
\item [2017 Fall Final] A rectangular fence consists of three sides costing \$2 per meter and one side costing \$1 per meter. If the area of the rectangle is 12 square meters, find the dimensions that minimize the cost of the fence.
\item [2017 Spring Exam 3] Let $p(x) = 100 - 2x$ be the price in dollars per cake a bakery can charge if it sells $x$ cakes. What cake price will maximize revenue?
\item [2017 Spring Exam 3] A rectangular open-topped aquarium is to have a square base and volume \SI{5}{m^3}. The material for the base costs \$10 per \si{m^2}, and the material for the sides costs \$1 per \si{m^2}. What dimensions minimize the cost of the aquarium?
\item [2017 Spring Final] A farmer has 20 feet of fencing and wants to fence off a rectangular region that borders a straight river. The farmer needs no fencing along the river. What dimensions will maximize the fenced-in area?
\item [2016 Fall Exam 3] A rectangular fence consists of three sides costing \$2 per meter and one side costing \$1 per meter. If the area of the rectangle is 12 square meters, find the dimensions that minimize the cost of the fence.
\item [2016 Fall Final] Find the dimensions of a cylinder with total surface area $6\pi$ square meters, including top and bottom, that maximizes its volume.
\item [2016 Spring Exam 3] A farmer has 24 feet of fencing and wants to fence off a rectangular area that borders a straight river. The farmer needs no fencing along the river. What dimensions will maximize the fenced-in area?
\item [2016 Spring Exam 3] A rectangular open-topped aquarium is to have a square base and volume \SI{8}{m^3}. The material for the base costs \$2 per \si{m^2}, and the material for the sides costs \$1 per \si{m^2}. What dimensions minimize the cost of the aquarium?
\item [2016 Spring Final] Suppose you want to enclose a \SI{25}{ft^2} rectangular area with fencing. What is the minimum length of fencing needed?
\item [2015 Fall Exam 3] A landscaper is designing a fence along the four sides of a rectangular garden, which is to have an area of 5000 square feet. The fencing for three sides costs \$10 per foot, but the fencing along the front side of the garden will cost \$30 per foot. Find the length and width of the garden in order to minimize the total cost.
\item [2015 Spring Exam 3] If a bakery charges $x$ dollars per cake, it makes a total profit of $P(x) = -x^2 + 100 x - 30$. If the bakery wants to maximize profit, what should it charge per cake?
\item [2015 Spring Exam 3] Find the dimensions of the box with square base that has volume 8 and minimal surface area.
\item [2014 Fall Exam 3] A homeowner with 16 feet of fencing wants to enclose a rectangular area against the side of her house. What dimensions will maximize the fenced-in area? (Note that 3 sides of the rectangle will be formed from fencing, and the house will serve as the fourth side of the rectangle.)
\item [2014 Fall Exam 3] A rectangular open-topped box is to have a square base and volume \SI{12}{ft^3}. If material for the base costs \$3 per \si{ft^2} and material for the sides costs \$1 per \si{ft^2}, what dimensions minimize the cost of the box?
\item [2014 Spring Exam 3] What is the smallest perimeter possible for a rectangle of area \SI{4}{ft^2}?
\end{enumerate}

\newpage
\section{Integrals}
I will separate the integral problems based on whether they require u-sub or not.
This is something one should be able to determine without being told.

\subsection{Basic (no u-sub)}
These should be doable by inspection. \\
Some may use ``guess and fudge'', i.e., simple u-sub (where $u = ax + b$ and $\dif u = a \dif x$, for $a,b\in\RR$). \\
Some use geometry / symmetry property of functions.

\begin{enumerate}
\item[2023 Spring Final \#1B] $\int\left( \frac{2}{t^3} + 2\sqrt{t} \right) \dif t$
\item[2022 Fall Exam 3 \#6A] $\int 2x+2x^{-2} + 3x^{2022} \dif x$
\item[2022 Fall Exam 3 \#6B] $\int \sqrt{x} + \frac{5+x^2}{1+x^2} \dif x$
\item[2022 Fall Exam 3 \#6C] $\int e^x + \frac1x + \sec x \tan x \dif x$
\item[2022 Spring Final \#1B] $\int \left( \frac1{t^2} - \sqrt t \right) \dif t$
\item[2022 Spring Final \#2A] $\int x\sqrt{x-1}\dif x$
\item[2021 Fall Final \#5A] $\int \frac{1}{x^2+4} + 3\sin x + \frac 1x \dif x$ \footnote{u-sub not actually needed if know general arctan integral formula}
\item[2021 Fall Final \#5B] $\int \frac{u^2+3}{\sqrt u} \dif u$
\item[2021 Fall Final \#6A] $\int \frac{(\ln x)^4}{x} \dif x$
\item [2019 Fall Exam 3 \#1C] $\int \left( \frac{1}{x^2} + 3\cos x \right) \dif x$
\item [2019 Fall Exam 3 \#1E] $\int_0^\pi (2\sin \theta+1) \dif \theta$
\item [2019 Fall Final \#1B] $\int \left(2\sqrt x - \frac{1}{x^2}\right) \dif x$
\item [2018 Fall Exam 3] $\int (\sqrt x + 6\sec^2(x) - 5) \dif x$
\item [2018 Fall Exam 3] $\int_0^4 (e^x - 3) \dif x$
\item [2018 Fall Final] $\int \left( \frac 2x - \sqrt{x} \right) \dif x$
\item [2017 Fall Exam 3] $\int \frac{\sqrt x - \sqrt 2 x^5}{x} \dif x$
\item [2017 Fall Final] $\int e^{5x} - \frac{1}{\sqrt{4-x^2}} \dif x$ \footnote{easier if you know general arcsine integral formula}
\item [2017 Spring Final] $\int (\sqrt x + \cos x - 5) \dif x$
\item [2016 Fall Final] $\int \sin(\pi x/ 2) + 2^x - \frac{1}{\sqrt{1-x^2}} \dif x$
\item [2017 Fall Exam 3] $\int_{-2}^2 \sin^3(5x) + \sqrt{4-x^2} \dif x$
\item [2016 Fall Exam 3] $\int \frac{x^2 - 7x}{x^3} \dif x$
\item [2016 Fall Exam 3] $\int_{-2}^2 \left( \frac{\sin x}{1+x^2} + \cos(\frac{\pi}4 x) \right) \dif x$
\item [2016 Spring Exam 3] $\int (\sec^2 x + 4) \dif x$
\item [2016 Spring Exam 3] $\int (\sqrt x + 5e^x) \dif x$
\item [2015 Fall Exam 3] $\int (x^2 + 4)^2 \dif x$
\item [2015 Fall Exam 3] $\int_1^2 2^t \dif t$
\item [2015 Spring Exam 3] Find the most general antiderivative of $\sec^2 x + 3x^4 + 2$.
\item [2014 Fall Exam 3] $\int \left( \cos x + 4x + \frac{1}{x} \right) \dif x$
\item [2014 Fall Exam 3] $\int \left( 3e^x + 4\sin x + 7\sec^2 x \right) \dif x$
\item [2014 Spring Exam 3] $\int (7 + 2x + 3e^x) \dif x$
\item [2014 Spring Exam 3] $\int (\sec^2 \theta + \cos \theta) \dif \theta$
\end{enumerate}

\subsection{U-sub}
\begin{enumerate}
\item[2023 Spring Final \#2A] $\int_3^4 x\sqrt{x-3} \dif x$
\item[2021 Fall Final \#5C] $\int \cos x \sin^3 x \dif x$
\item[2021 Fall Final \#6B] $\int 3t^2 \sqrt{1+t^3} \dif t$.
\item[2021 Fall Final \#6C] $\int_{-2}^{-1} x\sqrt{x+2} \dif x$
\item[2019 Fall Exam 3 \#7A] $\int_0^{\pi/2} e^{2\sin \theta} \cos\theta\dif\theta$
\item[2019 Fall Exam 3 \#7B] $\int x\sqrt{x-2}\dif x$
\item[2019 Fall Final \#2A] $\int \frac{\sec[2](2\theta)}{\tan[5](2\theta)} \dif \theta$
\item[2018 Fall Exam 3] $\int_0^{\pi/2} 2\sin^3(\theta) \cos(\theta) \dif \theta$
\item[2018 Fall Exam 3] $\int x \sqrt{5+x} \dif x$
\item[2018 Fall Final] $\int \frac{\cos(\ln x)}{x} \dif x$
\item[2017 Fall Exam 3] $\int 3x \sine(5x^2) \dif x$
\item[2017 Fall Exam 3] $\int x \sqrt{x+2} \dif x$
\item[2017 Fall Final] $\int_1^e \frac{(\ln x)^2}{x} \dif x$
\item[2017 Fall Final] $\int \sin^5(2x) \cos(2x) \dif x$
\item[2017 Spring Final] $\int \frac{\sqrt{\ln x}}{x} \dif x$
\item[2017 Spring Final] $\int_0^{\pi/2} \sin^4 \theta \cos \theta \dif \theta$
\item[2016 Fall Exam 3] $\int \sqrt{\tan x}\sec^2 x \dif x$
\item[2016 Fall Exam 3] $\int \frac{(\ln x)^3}{x} \dif x$
\item[2016 Fall Final] $\int \tan^3(2x) \sec^2(2x) \dif x$
\item[2016 Fall Final] $\int_0^1 \frac{x+2}{x^2+4x+1} \dif x$
\item[2016 Spring Final] $\int_0^{\pi/2} \sin^3 \theta \cos \theta \dif \theta$
\item[2016 Spring Final] $\int t \sqrt{t^2 + 3} \dif t$
\item[2015 Fall Exam 3] $\int \tan^3 x \sec^2 x \dif x$
\end{enumerate}

\subsection{with initial conditions (Initial Value Problem)}
\begin{enumerate}
\item [2023 Spring Final \#4] Find $f(x)$ if $f'(x) = \cos(x) + 2$ and $f(0) = 5$.
\item [2022 Spring Final \#4] Find $f(x)$ if $f'(x) = e^x+2$ and $f(0) = 7$.
\item [2019 Fall Exam 3 \#3] Find $f(x)$ if $f''(x) = e^x+2$, $f'(0) = 7$, and $f(0) = -5$.
\item [2019 Fall Exam 3 \#6] Suppose that a particle has position $s(t)$ feet at time $t$ seconds and a velocity function $s'(t) = 3t^2 - 2t$ ft/s. If $s(0) = 5\ \unit{ft}$, find $s(2)$.
\item [2018 Fall Exam 3] Find $f(x)$ if $f''(x) = 6x$, $f'(0) = 1$, and $f(0) = 2$.
\item [2017 Fall Exam 3] Solve the initial value problem for $f(t)$: $f'(t) = 2e^{-2t}$, $f(0) = 1$.
\item [2017 Fall Final] $f'(t) = 4t^3 - \sin t$, $f(0) = 1$.
\item [2016 Fall Exam 3] Solve the initial value problem for $f(t)$: $f'(t) = 4e^{3t}$, $f(0) = 5$.
\item [2016 Fall Final] $f'(t) = \sqrt{t}$, $f(1) = 2$.
\item [2016 Spring Exam 3] Find the function $v(x)$ satisfying $v''(x) = 2$, $v'(0) = -3$, and $v(0) = 5$.
\item [2015 Fall Exam 3] A ball thrown vertically from the roof of a building 150 feet tall hits the ground 3 seconds later. Was the ball thrown upward or downward? With what speed was it thrown? (Recall that acceleration due to gravity is $a = -32 \si{ft/sec^2}$.) \footnote{This is a neat problem.}
\item [2015 Spring Exam 3] Find the function $g(x)$ satisfying $g'(x) = \sin x + 1$ and $g(0) = 3$.
\item [2014 Fall Exam 3] Find $v(x)$ if $v''(x) = 6x+ 2$, $v'(0) = 1$, and $v(0) = 2$.
\item [2014 Spring Exam 3] Find the function $k(x)$ provided that $k'(x) = 2x^3 + 3x + 2$ and $k(0) = 2$.
\end{enumerate}

\subsection{Using linearity property}
\begin{enumerate}
\item [2015 Fall Exam 3] Suppose that $\int_0^6 f(x) \dif x = 9$, $\int_4^6 f(x) \dif x = 5$, and $\int_0^4 g(x) \dif x = 8$. Compute $\int_0^4 (5f(x) - 3g(x)) \dif x$.
\end{enumerate}

\section{FTC part 1}
\begin{enumerate}
\item[2023 Spring Final \#1C] $\dod{}{x} \int_x^2 \cos(e^t) \dif t$
\item[2022 Spring Final \#1C] $\dod{}{x} \int_x^7 \cos(\sin t)\dif t$
\item[2021 Fall Final \#3C] $\dod{}{x} \left( \int_0^{x^2} \cos(u+3) \dif u \right)$
\item [2019 Fall Exam 3 \#1D] $\dod{}{x} \int_x^5 3t^7 \cosine(t^9)\dif t$
\item [2019 Fall Final \#1C] $\dod{}{x} \int_x^7 \sine(t^5) \cosine(t^2+1) \dif t$
\item [2018 Fall Exam 3] $\dod{}{x} \int_3^x e^{2t} \sine(t^3) \dif t$
\item [2018 Fall Final] $\dod{}{x}\int_x^5 e^{\sin t} \dif t$
\item [2017 Fall Exam 3] $\dod{}{x} \int_2^x \frac{\cos(t^2)}{2+t} \dif t$
\item [2017 Fall Exam 3] $\dod{}{x} \int_{x^3}^5 \frac{\cos(t^2)}{2+t} \dif t$
\item [2017 Spring Final] $\dod{}{x} \int_3^x t \cdot \sine(t^3) \dif t$
\item [2016 Fall Exam 3] $\dod{}{x} \int_2^x \frac{\sin t}{1+t} \dif t$
\item [2016 Fall Exam 3] $\dod{}{x} \int_2^{x^3} \frac{\sin t}{1+t} \dif t$
\item [2016 Spring Final] $\dod{}{x} \int_0^x e^{\cos t} \dif t$
\item [2015 Fall Exam 3] Define $F(x) = \int_1^x \frac{\sin(\frac{\pi t}{6})}{t^2} \dif t$. Find an equation of the tangent line to $y = F(x)$ at $x = 1$. \footnote{interesting twist}
\end{enumerate}

\section{Approximate integral with Riemann sum (reading from graph)}
For some of these problems, a graph is given. 
Will not reproduce these.
\begin{enumerate}
\item [2023 Spring Final \#3]
\item [2022 Fall Exam 3 \#7A]
\item [2022 Spring Final \#3]
\item [2019 Fall Exam 3 \#4]
\item [2019 Fall Final \#12]
\item [2017 Fall Exam 3] Estimate the area below the curve $y = \sqrt x$ over the interval $[1, 4]$ using $L_3$. Sketch a graph of the curve and illustrate the rectangles used on the graph.
\item [2017 Fall Final] Estimate the area below the curve $y = x^2 + 2$ over the interval $[0, 3]$ using $R_3$. Make a sketch of the curve with the rectangles used.
\item [2016 Fall Exam 3] Estimate the area below the curve $y = \sqrt x + 1$ over the interval $[0, 6]$ using $R_3$. Sketch a graph of the curve and illustrate the rectangles used on the graph.
\item [2016 Fall Final] Estimate the area below the curve $y = x^2$ over the interval $[0, 2]$ using $L_4$. Make a sketch of the curve with the rectangles used.
\item [2015 Fall Exam 3] Approximate the area under the curve $y = 12x - 4x^2$ between $x = 0$ and $x = 2$ using four rectangles and the right endpoint method.
\item [2015 Spring Exam 3] Estimate the area between $y = x^2$ and the $x$-axis over the interval $[0, 4]$. Use $n = 2$ rectangles, taking the sampling points to be the midpoints (in other words, compute $M_2$). Sketch the rectangles on the graph.
\item [2014 Fall Exam 3] Estimate $\int_0^6 (x^2 + 1) \dif x$ by using $n = 3$ subintervals, taking the sampling points to be midpoints (in other words, compute $M_3$). Sketch the rectangles on a graph.
\end{enumerate}


\section{Definite integral as geometric area (reading from graph)}
Wont duplicate problem. Just list location.
Sometimes has FTC part 1 as part of the problem
\begin{enumerate}
\item [2023 Spring Final \#8]
\item [2022 Fall Exam 3 \#7B]
\item [2022 Spring Final \#8]
\item [2021 Fall Final \#11] This doubles as a first derivative interpretation and a net change theorem type of problem
\item [2019 Fall Exam 3 \#2]
\item [2019 Fall Final \#11] Part B is a FTC Part 1 problem
\end{enumerate}

\section{Net Change Theorem}
Include units.

\begin{enumerate}
\item [2023 Spring Final \#12] Suppose that a particle has position $s(t)$ feet at time $t$ seconds and a velocity function $s'(t) = t \cdot e^{-t^2}$ ft/sec. Find the displacement from time $t=0$ seconds to time $t=1$ second.
\item [2022 Spring Final \#12] Suppose that a particle has position $s(t)$ feet at time $t$ seconds and a velocity function $s'(t) = \sin^3(t)\cos(t)$ ft/sec. Find the displacement from time $t=0$ seconds to time $t=\frac\pi2$ seconds.
\item [2019 Fall Exam 3 \#5]
\item [2019 Fall Final \#5] Same as 2023 Spring Final \#12
\item [2018 Fall Exam 3] Suppose a particle has position $s(t)$ feet at time $t$ seconds and a velocity function $s'(t) = 3\cos(t)$ ft/s. Find the displacement from time $t = 0$ seconds to time $t = \pi/2$ seconds.
\item [2018 Fall Final] Suppose a particle has position $s(t)$ feet at time $t$ seconds and a velocity function $s'(t) = t \cdot \sine(\pi t^2)$ ft/s. Find the displacement from time $t=0$ to time $t = 1$.
\item [2016 Fall Exam 3] An object moves along a straight line with velocity $v(t) = 4 - t^2$ m/sec. Find the displacement of the object over the time interval $[0, 3]$ seconds. Find the total distance the object travels over the same time interval.
\item [2015 Fall Exam 3] An object moves along the $x$-axis with velocity $v = 12t^3 - 12t^2 $ \si{cm/sec}. Find the total distance traveled for the interval $-1 \le t \le 2$ seconds.
\end{enumerate}

\section{Area between curves (final)}
\begin{enumerate}
\item[2023 Spring Final \#5] Find the area between the curves $y=4$ and $y=x^2$.
\item[2022 Spring Final \#5] Find the area between the curves $y=x$ and $y=x^2$.
\item[2021 Fall Final \#10] Find the area of the region between the $y$-axis, the curve $y=6-x$ and the curve $y=x^2$
\item[2019 Fall Final \#8] Find the area between the curves $y=2x$ and $y=x^2$.
\item[2018 Fall Final] Same as 2023 Spring Final \#5 %Find the area between the curves $y = 4$ and $y = x^2$
\item[2017 Fall Final] Calculate the area of the region with $x\ge 0$ bounded by the $y$-axis, the parabola $y = x^2 - 2x$ and the line $y = 6-x$.
\item[2017 Spring Final] Same as 2023 Spring Final \#5 %Find the area bounded between $y = 4$ and $y = x^2$.
\item[2016 Fall Final] Calculate the area between $y = 8-x^2$ and $y = x+2$.
\item[2016 Spring Final] Find the area bounded between $y = 2x^2$ and $y = 3 - x^2$.
\end{enumerate}

\section{Volume of Solids of Revolution (final)}
\begin{enumerate}
\item[2023 Spring Final \#9] Find the volume of the solid obtained by rotating the region bounded by $y=\sqrt{\sin x}$ and $y=0$ between $x=0$ and $x=\pi$ around the $x$-axis.
\item[2022 Spring Final \#9] Find the volume of the solid obtained by rotating the region bounded by $y=x^3$, $y=0$, and $x=1$ around the $x$-axis.
\item[2019 Fall Final \#3] Same as 2023 Spring Final \#9 %Find the volume of the solid obtained by rotating the region bounded by $y=\sqrt{\sin x}$ and $y=0$ between $x=0$ and $x=\pi$ around the $x$-axis.
\item[2019 Fall Final \#4] Find the volume of the solid obtained by rotating the region bounded by $y=\frac1x$, $y=0$, $x=1$, and $x=2$ around the $y$-axis.
\item[2018 Fall Final] Rotate the region bounded by $y=x^3$, $x = 0$, and $y = 1$ around the $x$-axis. Find the volume.
\item[2018 Fall Final] Rotate the region bounded by $y=x^2$ and $y = x$ around the $y$-axis. Find the volume.
\item[2017 Fall Final] Rotate the region bounded by the $y$-axis, $y=x$, and $y=3 + \frac12 x$ about the $y$-axis. Set up the integral for the volume.
\item[2017 Fall Final] Rotate the region bounded by the $y$-axis, $y=x$, and $y=3 + \frac12 x$ about the $x$-axis. Set up the integral for the volume.
\item[2017 Spring Final] Find the volume of the solid obtained by rotating the region bounded by $y = x$ and $y = x^2$ around the $x$-axis.
\item[2016 Fall Final] Take the region bounded between the curves $y = 4x - x^3$ and $y = x^2$, with $x \ge 0$, and rotate around the $x$-axis. Set up the integral for the volume. (The curves don't intersect at a nice $x$-value.)
\item[2016 Fall Final] Take the region bounded between the curves $y = 4x - x^3$ and $y = x^2$, with $x \ge 0$, and rotate around the $y$-axis. Set up the integral for the volume. (The curves don't intersect at a nice $x$-value.)
\item[2016 Spring Final] Find the volume of the solid obtained by rotating the region bounded by $y = 4$ and $y = x^2$ around the $x$-axis.
\end{enumerate}

\end{document}
